\section{Discussion}
\label{sec:discussion}

\subsection{The core finding}

The rolling intermarket dependency between Bitcoin and conventional assets can be predicted out of sample. The result is not an artefact of in-sample overfitting: the signal withstands strict temporal-ordering constraints for all six pairs and four window lengths. Predictability depends almost entirely on persistence, since the rolling correlation shows strong serial dependence, and its current value is largely determined by its recent past.

These two observations are consistent with each other: the dependency is predictable mainly because it is persistent. The same pattern appears again and again over time, which makes the dependency easier to predict.

\subsection{Why simple models dominate}

The fact that Ridge, AR(1), and HAR finish at almost the same level at the top of the ranking says something about the data-generating process. A rolling Pearson correlation computed over $w$ days shares $w-1$ observations with the previous window, which mechanically creates strong autocorrelation. After the Fisher transformation, the first-order lag captures the vast majority of the forecastable variance. In these circumstances any reasonable model relies heavily on the most recent value, and a model that does little else is difficult to surpass.

That Ridge exceeds AR(1) by only 0.0003 RMSE (0.0656 versus 0.0659) is not a crucial advantage of machine learning; it reflects the fact that regularised linear regression, when scaled correctly, restores the same autoregressive signal with a slightly richer feature set. HAR corresponds almost exactly to AR(1), using three time scales (daily, weekly, and monthly) that match the overlapping-window memory structure of the target. The convergence of three structurally different methods to almost the same accuracy indicates that the prediction problem has an almost linear autoregressive structure, and that none of these methods extracts significantly different information.

The pattern is well documented in financial forecasting \cite{welch2008}: signal-to-noise ratios in financial markets are low, and the additional information in auxiliary features is often dominated by the autoregressive component. In this study the pattern became apparent only after a long hyperparameter search over tree configurations had been carried out.

The much lower performance of the tree methods is still useful evidence. It is diagnostic: the fact that well-configured XGBoost and Random Forest models cannot match Ridge or AR(1) indicates that the dependency process contains little exploitable nonlinear structure at the one-step horizon, at least within the scope of the feature set used here.

\subsection{ML compared with DCC-GARCH: a fairer comparison}

Ridge and ElasticNet have lower RMSE than DCC-GARCH in the reported experiments, and the descriptive gap is large: an average RMSE of 0.066 for Ridge compared with 0.214 for DCC-GARCH, more than a factor of three. The reported DM statistics are positive, but their zero-lag HAC standard errors do not account for autocorrelation induced by the overlapping rolling target; the associated $p$-values are therefore not used as formal significance evidence.

The correct interpretation is not that DCC-GARCH is a bad model. DCC-GARCH is a consistent conditional covariance model, and it produces instantaneous conditional correlation, whereas the goal here is a smoothed moving average. Projecting the DCC-GARCH output onto this moving target results in a systematic lag that the model was never intended to eliminate. Machine learning models trained directly on the Fisher-transformed moving target do not suffer from this mismatch, as they are optimized from the outset to achieve the same target.

A second distinct issue is data harmonization. Since conventional data series are fed into Bitcoin's seven-day calendar, the GARCH components for the benchmark are estimated based on inverted series in which approximately 30\% of the observations consist of forced zeros on weekends and holidays (Appendix A). This degrades the volatility estimates used in calculating the conditional correlation and further widens the gap with machine learning models. Re-estimating the DCC-GARCH model using only regular trading days would help identify this effect and is planned for future work. However, the robustness check for weekends in Appendix A shows that the constancy of the target itself—a property utilized by machine learning models—remains virtually unchanged under the calendar adjustment.

We should be honest about why there is a performance gap. If we simply say that machine learning is better than the econometric model without considering these differences, we make the results look stronger than they actually are.

\subsection{Strengths and limitations of the investor signal}

The signal layer gives moderate results in every dimension: AUC of approximately 0.48 to 0.56, balanced accuracy from below 0.50 to about 0.53 depending on the pair and window, and, for the logistic model, exit rates from 16 to 35\%. None of these numbers is strong in itself. Together they describe a signal that, on average, classifies stress regimes somewhat better than chance, activates with a reasonable frequency, and carries more information about equity stress than about precious-metal or dollar-index stress.

The equity result is amenable to economic interpretation. Bitcoin and the S\&P 500 share exposure to global risk appetite, funding liquidity, and sentiment cycles. When the dependency between them is rising or shifting, this often reflects broader market-regime transitions that also increase equity downside risk. The cryptocurrency feature set is reasonably well specified for that channel.

Gold and silver respond to a broader range of factors, including real interest rates, inflation expectations, geopolitical risks, and demand for safe-haven assets. Bitcoin's price movements do not reflect these factors very well. Thus, the weaker signal for these assets reinforces the same idea: the set of features captures some of the factors causing market stress, but not all of them.

In practice the signal belongs in the toolbox as an auxiliary risk indicator; used as a standalone alarm it would disappoint. Its natural place is next to fundamental valuation, macro analysis, and position-sizing rules that raise the bar for accepting risk when the dependency-based probability of a stress regime is elevated.

\subsection{Methodological choices that mattered}

Several design decisions were important. The most important was running DCC-GARCH through the same expanding-window walk-forward as every other model: a full-sample DCC fit would have given the benchmark access to future data and made the comparison invalid. The Fisher transformation mattered because the bounded Pearson correlation distorts errors close to $\pm 1$; the transformed target is closer to Gaussian, and RMSE becomes a sensible loss over the entire range. Reporting all four window lengths allowed the study to show honestly how much predictability depends on the autocorrelation structure, which the progression from $w=14$ to $w=90$ makes visible. Finally, reporting several metrics and paired DM diagnostics showed that DCC-GARCH has higher average error in every reported pair and window; formal inference requires a HAC specification that accounts for the rolling-window overlap.

\subsection{Statements this thesis does not make}

Several claims are explicitly not being made.

Bitcoin does not predict equity returns. It partially predicts the evolving dependency between Bitcoin and equities, which is a weaker and more precise statement. The signal layer estimates the probability of increased downside risk for conventional assets; it says nothing about the direction, size, or timing of the movement within a flagged period.

The results are based on predicting one step ahead. Multi-step performance may differ significantly, especially for machine learning, where the autoregressive advantage weakens as the horizon expands. This is an open question.

These results are based on one-step-ahead forecasting. Multi-step performance may vary significantly, especially in the case of machine learning, where the advantage of autoregression diminishes as the forecasting horizon increases. This question remains open.

\subsection{Iterative model development: what changed across versions}
\label{subsec:iterations}

The final model presented in this thesis is the result of several months of iterative refinement. This section documents that process, because the sequence of changes is itself an empirical result about the nature of the forecasting problem.

\paragraph{Version 1 (September 2025): minimal baseline.}
The first working prototype had seven models: the last-value baseline, AR(1), Ridge, ElasticNet, Random Forest, Gradient Boosting, and XGBoost. Feature engineering was very simple. We used four lags of the rolling correlation, added short-window volatilities for both assets, and included lagged returns up to lag 5. We did not include a DCC-GARCH benchmark at this stage.

Hyperparameter tuning was done ad hoc. Ridge used alpha $= 1.0$, ElasticNet used alpha $= 0.01$, and XGBoost used 1000 estimators with a learning rate of 0.05. We did not scale the features for the linear models.

For BTC-USD versus S\&P 500 at $w = 30$, the best out-of-sample RMSE was 0.0897 for Ridge. The AR(1) model performed comparably.

\paragraph{Version 2 (November 2025): DCC-GARCH and signal layer.}
The second major milestone introduced the DCC-GARCH econometric benchmark via a walk-forward protocol that matched the expanding-window setup used for machine-learning models. This was a critical correctness fix: an earlier full-sample DCC fit would have granted the benchmark access to future data, invalidating the comparison. The investor signal layer was added as a second stage, classifying whether conventional assets were entering stress regimes using the rolling-dependency features. Core model hyperparameters and feature sets were unchanged from Version 1.

\paragraph{Version 3 (February 2026): stabilisation.}
In this version, the pipeline was stabilized as a reliability baseline. There were no improvements in metrics; the primary focus was on reliability. We added a fallback version of XGBoost from GPU to CPU, improved exception handling, and made metric export consistent. Out-of-sample results were nearly identical to those of Version 2, confirming that the performance degradation in Version 1 was not caused by implementation errors.

\paragraph{Version 4 (May 2026): feature expansion, regularisation revision, and HAR.} The final version incorporated several targeted improvements informed by the Version 1--3 plateau:

\begin{itemize}
\item Feature engineering expanded from $\approx$15 to $\approx$35 features. New features include momentum and z-score of the rolling correlation (dep\_mom\_5, dep\_mom\_20, dep\_zscore), extended lags (dep\_lag20, dep\_lag60), a 60-day volatility window, the volatility ratio between assets, squared-return proxies, and a short-run correlation momentum term. The Heterogeneous AutoRegressive (HAR) model \cite{corsi2009har} was introduced as a theoretically grounded alternative to AR(1). HAR uses three temporal scales simultaneously, fitted on the Fisher-$z$ target exactly as in Eq. (7): $\hat{z}_{t+1} = \phi_0 + \phi_d\,z_t + \phi_w\,\bar{z}^{(5)}_t + \phi_m\,\bar{z}^{(22)}_t + \varepsilon_{t+1}$, where $\bar{z}^{(k)}_t$ denotes the equally-weighted average of the previous $k$ observations. Originally developed for realised variance, HAR is well-suited to rolling correlations because the overlapping-window structure creates natural heterogeneous memory at daily, weekly, and monthly horizons. Ridge's penalty strength (scikit-learn's alpha, the $\lambda$ of Eq. (8)) was reduced from alpha $= 1.0$ to alpha $= 0.5$ and an explicit StandardScaler was added. The combined effect was dramatic: Ridge RMSE at BTC-USD vs S\&P 500 $w = 30$ improved from 0.0897 (Version 1) to 0.0651, nearly matching AR(1). The cross-validated optimum actually sits even lower (around alpha $= 0.001$; see Figure~\ref{fig:app_ridge_path}), but the CV curve is essentially flat across this range, so the slightly firmer alpha $= 0.5$ was kept for a more stable fit at negligible cost in RMSE. Conversely, XGBoost regularisation was tightened (learning rate 0.05 $\to$ 0.02, added $L_1/L_2$ penalties), reflecting the insight that the expanded feature set increased overfitting risk for tree-based methods. The walk-forward ensemble now weights each constituent model by the inverse of its RMSE over the most recent 60 out-of-sample predictions, instead of equal weights. This allows the ensemble to adapt dynamically to regime shifts in which different models temporarily dominate.
    \item Bootstrap confidence intervals and sensitivity analyses. 500-replicate bootstrap CIs were added for RMSE and $R^2$, and a refit-frequency sensitivity sweep across $\{5, 10, 20, 42, 63\}$ days (which includes the default of 20) was added to verify that $refit_every = 20$ is not a cherry-picked choice.
\end{itemize}

\paragraph{The lesson of the four versions.}
The biggest lesson from the four model versions is that adding more variables did not improve the tree-based models. Increasing the feature set from 15 to 35 variables did not make XGBoost better. Its RMSE stayed almost the same or became slightly worse after the expansion. On the other hand, Ridge improved substantially because it uses the full feature set in a linear way with regularisation. This result fits with the main point of the thesis. The forecasting target has a structure that is almost linear and autoregressive. The additional information from volatility ratios, momentum, and return squares is useful for the regularised linear model, but it does not give gradient-boosted trees useful nonlinear patterns to exploit.

For future work, this means that even more sophisticated features are unlikely to bridge the gap that still exists between AR(1) and machine learning in solving this problem. Real progress will most likely be achieved through a different approach to modeling, such as the use of multi-stage horizons, high-frequency data, or the simultaneous modeling of multiple correlation pairs.

% ─────────────────────────────────────────────────────────────────────────────
\subsection{Limitations}

The main limitations are: daily frequency (no intraday), one-step horizon, Pearson correlation as the only dependency measure, a single base cryptocurrency (Bitcoin), and no portfolio backtest. Each represents a natural extension that could sharpen or complicate the findings. None of them invalidates the core empirical results within the scope of what was actually tested.
